\documentclass[11pt]{article}

\title{NRR-Guarantee: Bounded Commit/Defer Guarantees Under Explicit Conditions}
\author{Kei Saito}
\date{March 20, 2026}

\begin{document}
\maketitle

\begin{abstract}
This manuscript initializes the current \texttt{Guarantee}-line paper surface for
the Non-Resolution Reasoning (NRR) program. The current version fixes the section
structure, package boundary, and evidence-import surface required to close the next
downstream paper under explicit provider, family, horizon, and observability
conditions. It now also fills the first bounded claim unit for the frozen early
ambiguity-turn surface. This remains a narrow entry claim rather than a final
multi-surface guarantee closure.
\end{abstract}

\section{Introduction}

NRR is not an anti-LLM framework. NRR does not replace standard LLM use. NRR
optimizes when to commit and when to defer, under explicit conditions.

The current downstream architecture is \textit{Guarantee-centered}: \texttt{Policy},
\texttt{Dynamics}, and \texttt{Policy-Verification} are supporting layers for the
\texttt{Guarantee} line rather than independent completion targets. Accordingly, the
purpose of this manuscript is to close a bounded downstream paper that absorbs those
layers as constrained inputs.

The present version is intentionally narrow. It fixes the manuscript surface that
will later host the bounded claim, but it does not yet promote any unverified
provider-universal, family-universal, or long-horizon statement.

\paragraph{Series Consistency Statement:}
Unsupported interpretations are never injected at a fixed evidence state,
hypothesis-set expansion is evidence-gated, and unobserved possibilities remain
explicitly unresolved. In this paper, the evidence-source clause is restricted to
the frozen downstream support layers and the directly imported observability
surfaces named for the current \texttt{Guarantee} package.

\paragraph{Series Extensibility Statement:}
NRR is a derivation-and-evaluation framework, not a closed recipe. The currently
evaluated policy/operator set is treated as an evaluated set under tested
conditions, not as an exhaustive set.

\section{Program Position And Scope}

This paper sits at the end of the current downstream sequence:

\begin{center}
\texttt{Energy -> Policy -> Dynamics + Policy-Verification -> Guarantee}
\end{center}

\noindent The role split is fixed as follows.

\begin{itemize}
\item \texttt{Energy}: measurement and bounded upstream handoff.
\item \texttt{Policy}: minimal state/action specification.
\item \texttt{Dynamics}: hold/degrade/break vocabulary across later checkpoints.
\item \texttt{Policy-Verification}: prospective matched-condition validation.
\item \texttt{Guarantee}: bounded claim under explicit conditions.
\end{itemize}

This manuscript must therefore do two things at once:
\begin{enumerate}
\item absorb the already fixed supporting layers without reopening their ownership
boundaries, and
\item state only those guarantee-shaped claims that remain supported after explicit
provider, family, horizon, observability, and failure-boundary restriction.
\end{enumerate}

\section{Supporting-Layer Inputs}

\subsection{Policy Input}

The current policy layer freezes the downstream state vocabulary
(\texttt{active\_objective}, \texttt{suspended\_objective},
\texttt{accepted\_direction}, \texttt{unresolved\_relation},
\texttt{quarantine}, \texttt{local\_relation\_block}) and the minimal action set
used to govern commit/defer behavior. The \texttt{Guarantee} line inherits this
vocabulary and must not redesign it locally.

\subsection{Dynamics Input}

The current dynamics layer contributes mechanism vocabulary only. In particular, it
provides the motifs \textit{Stable Resolution}, \textit{Clean Quarantine Hold},
\textit{Useful Continuation Hold}, \textit{Relation-Preservation Hold},
\textit{Premature Relation Collapse}, \textit{Repeated-Starter Attractor},
\textit{Generic Stall}, and \textit{Quarantine Leak}. These motifs describe where
policy-frozen behavior holds or breaks across later checkpoints, but they do not
yet constitute guarantee language.

\subsection{Policy-Verification Input}

The current policy-verification layer fixes the protocol rules that any prospective
\texttt{baseline vs Energy policy} comparison must satisfy before guarantee text is
allowed. This includes matched-condition comparison, pre-frozen claim units, narrow
family selection, explicit horizon bands, and main-text reporting of both gains and
failure boundaries.

\section{Claim Unit And Boundary Discipline}

Any later guarantee statement in this paper must fill a single bounded claim unit
with the following fields:

\begin{enumerate}
\item provider scope
\item family scope
\item horizon scope
\item observability scope
\item comparison scope
\item supported behavior
\item explicit failure boundary
\item unsupported extension
\end{enumerate}

No guarantee statement is valid if any field is omitted. This paper therefore
inherits the current schema rule that the claim surface must be chosen before
outcome inspection and must remain narrow enough to prevent provider-universal,
family-universal, or long-horizon overreach.

\subsection{First Filled Claim Unit}

The first explicit claim unit in the current manuscript is:

\begin{itemize}
\item provider scope:
  the tested provider pair only (\texttt{claude-sonnet-4-6} and
  \texttt{gemini-2.0-flash})
\item family scope:
  the frozen \texttt{priority-resolution} ambiguity-turn family only
\item horizon scope:
  the early decision boundary only
\item observability scope:
  \texttt{resolution\_attempt\_present},
  \texttt{unilateral\_commit\_absent}, and \texttt{appropriate}
\item comparison scope:
  matched baseline vs Energy-policy under the frozen early
  ambiguity-turn pilot conditions
\item supported behavior:
  non-unilateral priority resolution
\item explicit failure boundary:
  no downstream acceptance, side-topic/quarantine, or later-horizon closure is
  established by this fill
\item unsupported extension:
  no provider-universal, long-horizon, or broad downstream-quality claim is
  authorized
\end{itemize}

In manuscript form, the bounded claim is:

\begin{quote}
Under the tested provider pair, for the frozen \texttt{priority-resolution}
ambiguity-turn family, at the early decision boundary, and on the directly scored
fields \texttt{resolution\_attempt\_present},
\texttt{unilateral\_commit\_absent}, and \texttt{appropriate}, the Energy-policy
arm is supported relative to matched baseline only for non-unilateral priority
resolution. This support excludes downstream acceptance, side-topic/quarantine,
and later-horizon behavior, and it does not extend to provider-universal or broad
downstream-quality claims.
\end{quote}

\section{Evidence Import Surface}

The current evidence imports to be promoted into manuscript-facing text are:

\begin{itemize}
\item Program design:
  \texttt{energy\_to\_guarantee\_program\_design\_brief\_v1\_2026-03-17.md}
\item Policy:
  \texttt{policy\_minimal\_spec\_current\_v1\_2026-03-18.md}
\item Dynamics:
  \texttt{energy\_dynamics\_mechanism\_skeleton\_candidate\_v2\_2026-03-17.md}
  and
  \texttt{non\_evaluative\_principle\_dynamics\_closure\_note\_v1\_2026-03-19.md}
\item Policy-Verification:
  \texttt{policy\_verification\_pilot\_protocol\_current\_v1\_2026-03-18.md}
\item Guarantee schema:
  \texttt{energy\_guarantee\_claim\_schema\_candidate\_v2\_2026-03-17.md}
\end{itemize}

The current manuscript-facing task is not to widen these inputs, but to integrate
them into one paper whose claim boundary remains explicit and whose package
boundary can be externally reviewed.

\subsection{First Evidence Import}

The strongest current entry point for the \texttt{Guarantee} line is the frozen
early ambiguity-turn pilot. Under matched baseline vs Energy-policy conditions,
both tested provider slices moved from \texttt{0 / 24} to \texttt{24 / 24} on the
pilot's \texttt{appropriate} field, while baseline failure modes remained visible
as either ambiguity-preserving but unresolved responses or
unilateral/content-before-resolution behavior. The externally cross-checked
carry-forward read is therefore that the Energy-policy surface supports
non-unilateral priority resolution at the early decision boundary on the tested
provider pair. This is a bounded early-control claim only; it does not by itself
establish downstream followthrough, quarantine behavior, later-horizon stability,
or provider-universal validity.

\section{Reproducibility And Package Boundary}

The repository surface for the current line is intentionally minimal:
\begin{itemize}
\item a current TeX source
\item a current PDF snapshot
\item a working drafting outline
\item a supporting-layer evidence map
\item stable build and verification wrappers
\end{itemize}

This package is sufficient for writing and review-surface stabilization. It is not
yet sufficient for final external closure because the final bounded claim text,
explicit evidence imports, and external review artifacts still remain to be
completed.

\section{Discussion}

The non-evaluative principle should enter this line as bounded engineering
discipline rather than theorem inflation. In the current read, the practical issue
is not simply whether unresolved content exists, but whether later interaction
forces premature relation collapse or unnecessary comparative evaluation pressure
while the interaction is still in a retained state. The \texttt{Guarantee} paper is
the correct place to express that read as a bounded downstream claim surface.

\section{Current Completion Checklist}

The remaining manuscript/package tasks are:
\begin{enumerate}
\item add the next bounded claim units, if they remain supportable after the same
field-by-field restriction
\item expand manuscript-facing evidence imports from the current support memos
\item add explicit gains, regressions, and failure-boundary reporting across the
selected supported surfaces
\item pass independent Codex and Claude external review
\item close the git-ready package surface
\end{enumerate}

\section{Conclusion}

The current version establishes the initial reviewable package surface for
\texttt{Guarantee}. It fixes the section structure, inherits the current supporting
layer boundaries, and now adds a first explicit bounded claim unit for the early
ambiguity-turn surface while reserving broader downstream closure for later,
explicitly bounded revisions.

\end{document}
